\section{The exact block density}\label{sec:rg-step}

The block transformation has three factors: a new quadratic
fermion density, a determinant quotient, and a finite integral
over photon and fermion fluctuations. We derive this identity
before expanding its last factor. The integration at level
zero includes every chamber component and uses the chamber
Berezin measure. The algebraic identities in this section follow
\cite[\S13.1--13.3]{QED3Master}.

\subsection{Integration spaces and real backgrounds}
\label{subsec:rg-step-real-background}

Fix $k\in\Z_{\ge0}$, $\eta_k=L^{-k}$, odd $L\ge5$, and an
integer power $M$ of $L$ with $M\ge L^2$, with $\ell/(LM)$ odd and at least $13$.
The regulator parameter satisfies $0<\rho<1$.
Keep $b=1$ as a letter, and write
$b_L=b/L$ and
$b_k=b(1-L^{-1})/(1-L^{-k})$ for $k\ge1$.
Write $\mathscr V_k=\Sspace(1,\ell)$ for $k\ge1$, with
four spin components and counting measure, and
$\mathscr V_0=\Sspace^{\mathrm{ch}}(1,\ell)$, with
$192$ components per photon site and measure $1/48$.
The one-step averaging operator into $\Sspace(L,\ell)$ is
\[
 \mathcal A_k(B)=
 \begin{cases}
 Q(B),&k\ge1,\\
 Q(B)J^*,&k=0,
 \end{cases}
 \qquad
 \mathcal A_k^\sharp(B)=\mathcal A_k(-B)^T.
\]
The target measure is $L^3$. On these spaces set
\begin{equation}\label{eq:rg-step-quadratic-operators}
 \mathcal D_k(B)=
 \begin{cases}D_k(B),&k\ge1,\\D_B+\mu,&k=0,\end{cases}
 \qquad
 K_k(B)=\mathcal D_k(B)+b_L\mathcal A_k^\sharp(B)\mathcal A_k(B).
\end{equation}
Thus $K_k=D_k+b_LP_L$ for $k\ge1$ and
$K_0=F_0=D_B+\mu+b_LJP_LJ^*$.
We write $D_0(B)=D_B+\mu$, so that
$\mathcal D_k(B)=D_k(B)$ at every level.
When it exists, $\Gamma_k=K_k^{-1}$, and
$\delta Z_k(B)=\det K_k(B)$.
The transpose refers to the displayed bilinear measures.
For a real $B$, the phases have modulus one and
$\mathcal A_k^\sharp=\mathcal A_k^*$.

The normalized averages give
$Q(B)Q^T(-B)=1$. On a coarse block $B(z')$ the product
is
$L^{-3}\sum_{z\in B(z')}U_B(z',z)U_B(z',z)^{-1}=1$.
At level zero the additional cluster mean contributes
$48/48=1$. Similarly
$\mathcal Q_k\mathcal Q_k^\sharp=1$ for $k\ge1$.
In a unit block the photon count gives
$\eta_k^3\eta_k^{-3}=1$ and the cluster count again
gives $48/48=1$.
For real $B$ the inverse phases in the embeddings are
the complex conjugates of the averaging phases.
The products $\Pi_1=\mathcal Q_k^\sharp\mathcal Q_k$
and $\Pi_L=\mathcal Q_k^\sharp P_L\mathcal Q_k$ are
therefore orthogonal projections; so are $P_L$ and,
at level zero, $JP_LJ^*$. Their scalar spin kernels
commute with every constant spin matrix.

The fine resolvents used below are those of
\cite[\S4.2]{FermionBlockManuscript}. Explicitly, for $k\ge1$ put
\[
 F_\infty(B)=D_B+\mu+b_k\mathcal Q_k^\sharp\mathcal Q_k,
 \quad
 \mathcal B_{k,y}=(b_k+i\gamma_5y)^{-1}(1-P_L)
                  +(b_k+b_L+i\gamma_5y)^{-1}P_L,
\]
and define
\[
 F_y=F_\infty-b_k^2\mathcal Q_k^\sharp\mathcal B_{k,y}\mathcal Q_k,
 \qquad S_{k,y}=F_y^{-1}
\]
where the inverse exists. In particular $S_k=F_\infty^{-1}$.
The reconstruction maps are
$T_k(B)=\mathcal Q_k(B)(b_k^{-1}D_B+1+\mu/b_k)$ and
$\bar T_k(B)=\mathsf C^{-1}T_k(-B)\mathsf C$;
their recovery identities are \cite[Theorem~B(iv)]{NormalizedContraction}.
At zero these reconstruction maps are identities on chamber fields.

\begin{proposition}\label{prop:rg-step-real-inverses}
Let $B$ be real and $\mu>0$. The operators
$D_B+\mu$, $\mathcal D_k(B)$ and $K_k(B)$ are
invertible at every level at which they are defined.
For $k\ge1$,
\begin{equation}\label{eq:rg-step-real-coarse-inverse}
 D_k(B)^{-1}=b_k^{-1}
           +\mathcal Q_k(B)(D_B+\mu)^{-1}\mathcal Q_k^\sharp(B),
 \qquad \operatorname{Re}D_k(B)^{-1}\ge b_k^{-1}.
\end{equation}
Moreover $\operatorname{Re}D_k>0$ and
$\operatorname{Re}K_k>0$ as quadratic forms for
$k\ge1$, whereas $\operatorname{Re}K_0\ge\mu$.
Each $K_k$ satisfies
$K_k^*=\gamma_5K_k\gamma_5$.
For every real $y$, $K_k+i\gamma_5y$ is invertible,
and for $y\ne0$,
\begin{equation}\label{eq:rg-step-real-resolvent}
 \|(K_k+i\gamma_5y)^{-1}\|_{2\to2}\le |y|^{-1}.
\end{equation}
For $k\ge1$ the fine inverse $S_{k,y}$ also
exists for every real $y$.
In particular $\det D_k(0)$ and $\delta Z_k(0)$
are nonzero, including the determinants on
$\mathscr V_0$.
\end{proposition}

\begin{proof}
At a real connection the chamber hops are self-adjoint
unitary involutions. In the proof of
\exteq{6}{FermionBlockManuscript}, reindexing $x$ by $x^{(i)}$
uses
$U_B(x^{(i)},x)=\overline{U_B(x,x^{(i)})}$.
It gives
\[
 (f,\mathcal L_Bf)=\frac12\sum_iw_i
                  \int\dd x\,|R_i(B)f(x)-f(x)|_2^2\ge0.
\]
The direction vector changes sign under the same
involution; hence
$\sum_iw_i(\gamma\cdot\widehat d_i)R_i(B)$
is anti-self-adjoint. The cancellation
$\sum_iw_i\widehat d_i=0$ yields
$D_B^*=\gamma_5D_B\gamma_5$ and
$\operatorname{Re}D_B=(\rho/(2\eta_k))\mathcal L_B\ge0$.
Thus $\operatorname{Re}(D_B+\mu)\ge\mu$.
For $k\ge1$ adding the positive penalty
$b_k\Pi_1$ gives $\operatorname{Re}F_\infty\ge\mu$.
These inequalities imply injectivity and then
invertibility in the finite-dimensional spaces.

Put $G=D_B+\mu$ and $S_k=(G+b_k\mathcal Q_k^\sharp
\mathcal Q_k)^{-1}$. Multiplication of
$b_k^{-1}+\mathcal Q_kG^{-1}\mathcal Q_k^\sharp$
by $D_k=b_k-b_k^2\mathcal Q_kS_k\mathcal Q_k^\sharp$
gives
\[
 1+b_k\mathcal Q_kG^{-1}
      [1-(G+b_k\mathcal Q_k^\sharp\mathcal Q_k)S_k]
                                               \mathcal Q_k^\sharp=1.
\]
This proves \eqref{eq:rg-step-real-coarse-inverse}.
For an invertible operator $O$,
$\operatorname{Re}(u,Ou)=
\operatorname{Re}(Ou,O^{-1}Ou)$.
In particular $\operatorname{Re}O\ge0$ implies
$\operatorname{Re}O^{-1}\ge0$.
As $\mathcal Q_k^\sharp=\mathcal Q_k^*$ here,
the second term of $D_k^{-1}$ has nonnegative real
part. For a nonzero $v$, consequently,
\[
 \operatorname{Re}(v,D_kv)
       =\operatorname{Re}(D_kv,D_k^{-1}D_kv)
                           \ge b_k^{-1}\|D_kv\|_2^2>0.
\]
Adding $b_LP_L$ preserves strict positivity.
At $k=0$ the same conclusion follows directly by
adding $b_LJP_LJ^*$ to $D_B+\mu$.

The projections commute with $\gamma_5$.
It follows first that
$F_\infty^*=\gamma_5F_\infty\gamma_5$ and
$S_k^*=\gamma_5S_k\gamma_5$.
Taking the adjoint in the formula for $D_k$ and using
$\mathcal Q_k^\sharp=\mathcal Q_k^*$ gives
$D_k^*=\gamma_5D_k\gamma_5$. The penalty has the
same symmetry, proving it for $K_k$.
For $k=0$ this is already the symmetry of its three
terms. Hence
\[
 (K_k+i\gamma_5y)^*(K_k+i\gamma_5y)
                                      =K_k^*K_k+y^2,
 \qquad
 \|(K_k+i\gamma_5y)f\|_2^2
                         =\|K_kf\|_2^2+y^2\|f\|_2^2.
\]
For $y\ne0$ this proves injectivity and the inverse
bound; $y=0$ is covered by strict positivity.
For the fine inverse at $k\ge1$, use the block
matrix with diagonal blocks
$b_k+b_LP_L+i\gamma_5y$ and $F_\infty$,
and off-diagonal blocks
$-b_k\mathcal Q_k$ and $-b_k\mathcal Q_k^\sharp$.
The first diagonal block is invertible by its
idempotent decomposition
\exteq{53}{FermionBlockManuscript}.
Its two Schur complements are $F_y$ and
$K_k+i\gamma_5y$. The latter is invertible,
so the block matrix and then $F_y$ are invertible.
This proves existence of $S_{k,y}$.
The asserted normalizing determinants are therefore
nonzero at the real zero connection.
\end{proof}

\subsection{The next quadratic density}
\label{subsec:rg-step-scale-identification}

Assume now that $\ell/L$ is also admissible, and let
$B=B'_L$. At the outgoing level the mass is $L\mu$ and
the charge is $e_{k+1}=L^{1/2}e_k$.
The following identities hold whenever $S_k(B)$ and
$\Gamma_k(B)$ exist for $k\ge1$, or $\Gamma_0(B)$
exists for $k=0$. This invertibility domain is open
and invariant under $B\mapsto-B$: all elementary
entries are entire functions of the bond values,
and the transpose identities are
\begin{equation}\label{eq:rg-step-transpose}
 \mathcal D_k(B)^T=\mathsf C^{-1}\mathcal D_k(-B)\mathsf C,\qquad
 K_k(B)^T=\mathsf C^{-1}K_k(-B)\mathsf C,\qquad
 \Gamma_k(B)^T=\mathsf C^{-1}\Gamma_k(-B)\mathsf C.
\end{equation}
For the first identity reindex each hopping term by
$x^{(i)}$, using
$M_i(x^{(i)})^T=\mathsf C^{-1}M_i(x)\mathsf C$
and the reverse phase at $-B$. Transposition of the
averages reverses their connection in the product
$P_L$, and the scalar kernels commute with
$\mathsf C$. The remaining identities follow by
products and inversion. Invertible matrices have
analytic inverses on their open determinant domains.

Define the fine-to-$L$ average and its embedding by
\[
 \mathcal Q_k^+(B)=
 \begin{cases}Q(B)\mathcal Q_k(B),&k\ge1,\\Q(B)J^*,&k=0,\end{cases}
 \qquad
 \mathcal Q_k^{+,\sharp}(B)=\mathcal Q_k^+(-B)^T.
\]
They obey
$\mathcal Q_k^+\mathcal Q_k^{+,\sharp}=1$.
Set $\bar b_k=b_{k+1}/L$ and
\begin{equation}\label{eq:rg-step-fine-plus}
 F_k^+(B)=D_B+\mu+\bar b_k\mathcal Q_k^{+,\sharp}\mathcal Q_k^+,
 \qquad \Sigma_k^+(B)=F_k^+(B)^{-1}.
\end{equation}
For $k=0$ this is $F_0$ and its inverse $\Gamma_0$.
For $k\ge1$ it is the fine resolvent operator at
$y=0$, since
$\bar b_k=b_kb_L/(b_k+b_L)$.

To verify existence and the identities needed below,
write $\mathcal B_k=(b_k+b_LP_L)^{-1}$ for $k\ge1$
and consider
\[
 \begin{pmatrix}
 b_k+b_LP_L&-b_k\mathcal Q_k\\
 -b_k\mathcal Q_k^\sharp&F_\infty
 \end{pmatrix}.
\]
Its lower-right Schur complement is $K_k$.
Since $P_L$ is idempotent and scalar in spin,
$\mathcal B_k=b_k^{-1}(1-P_L)+(b_k+b_L)^{-1}P_L$.
The other Schur complement is
\[
 F_\infty-b_k^2\mathcal Q_k^\sharp\mathcal B_k\mathcal Q_k
    =D_B+\mu+\frac{b_kb_L}{b_k+b_L}
                                  \mathcal Q_k^\sharp P_L\mathcal Q_k
    =F_k^+.
\]
The two assumed inverses make the block matrix
invertible, and hence $F_k^+$ is invertible.
Equating the upper-left and lower-left blocks of
its inverse gives
\begin{equation}\label{eq:rg-step-block-inverse-identities}
 \Gamma_k=\mathcal B_k+
       b_k^2\mathcal B_k\mathcal Q_k\Sigma_k^+
                                     \mathcal Q_k^\sharp\mathcal B_k,
 \qquad
 H_k\Gamma_k=b_k\Sigma_k^+\mathcal Q_k^\sharp\mathcal B_k.
\end{equation}
This is the same Schur calculation as
Proposition~\extref{4.2}{FermionBlockManuscript}, retaining
the off-diagonal block for the mean field.

Let
\begin{equation}\label{eq:rg-step-unscaled-coarse}
 D_{k+1}^0(B)=b_L-
                       b_L^2\mathcal A_k(B)\Gamma_k(B)\mathcal A_k^\sharp(B).
\end{equation}
For $k\ge1$, $QP_L=Q$ and $P_LQ^T(-B)=Q^T(-B)$
give
$Q\mathcal B_k=(b_k+b_L)^{-1}Q$ and
$\mathcal B_kQ^T(-B)=(b_k+b_L)^{-1}Q^T(-B)$.
Insert these in \eqref{eq:rg-step-block-inverse-identities}
to obtain
\begin{equation}\label{eq:rg-step-compressed-identities}
 \begin{split}
 D_{k+1}^0(B)
       &=\bar b_k-\bar b_k^2
                            \mathcal Q_k^+\Sigma_k^+\mathcal Q_k^{+,\sharp},\\
 b_LH_k(B)\Gamma_k(B)\mathcal A_k^\sharp(B)
       &=\bar b_k\Sigma_k^+\mathcal Q_k^{+,\sharp}.
 \end{split}
\end{equation}
For $k=0$ both formulas follow directly from
$\bar b_0=b_L$, $H_0=1$,
$\Sigma_0^+=\Gamma_0$, and
$\mathcal Q_0^+=\mathcal A_0$.

Let $\mathsf R_L$ denote the scaling
$f_L(c,g)=L^{-1}f(c/L,g)$ on chamber spinors and
$v_L(z')=L^{-1}v(z'/L)$ on unit spinors.
These are bijections between the corresponding
incoming and outgoing spaces. Scaling preserves
the label and the nested block centers, and
$e_k\eta_kL^{-1/2}=e_{k+1}\eta_{k+1}$ preserves
each transport phase. Consequently
\[
 (D_B+\mu)\mathsf R_L
       =L^{-1}\mathsf R_L(D_{B'}+L\mu),\qquad
 \mathcal Q_k^+(B)\mathsf R_L
       =\mathsf R_L\mathcal Q_{k+1}(B'),
\]
and the corresponding embedding identity holds.
For the average, its $L^{-3}$ factor cancels the
$L^3$ change of the chamber index measure.
Thus
\begin{equation}\label{eq:rg-step-fine-scaling}
 F_k^+(B)\mathsf R_L
                   =L^{-1}\mathsf R_LF_\infty(B'),\qquad
 \Sigma_k^+(B)\mathsf R_L
                         =L\mathsf R_LS_{k+1}(B').
\end{equation}
In particular the outgoing inverse exists.
Equations \eqref{eq:rg-step-compressed-identities} and
\eqref{eq:rg-step-fine-scaling} now imply
\begin{equation}\label{eq:rg-step-coarse-scaling}
 D_{k+1}^0(B)\mathsf R_L
                       =L^{-1}\mathsf R_LD_{k+1}(B'),\qquad
 b_LH_k(B)\Gamma_k(B)\mathcal A_k^\sharp(B)\mathsf R_L
                       =\mathsf R_LH_{k+1}(B').
\end{equation}
All expressions in the second identity act on the
outgoing four-component field.

For independent outgoing Grassmann fields define
\begin{equation}\label{eq:rg-step-mean-fields}
 \Psi^c=b_L\Gamma_k(B)\mathcal A_k^\sharp(B)\Psi'_L,\qquad
 \bar\Psi^c=b_L\Gamma_k(B)^T\mathcal A_k(B)^T\bar\Psi'_L.
\end{equation}
Then, with $\psi'=H_{k+1}(B')\Psi'$ and
$\bar\psi'=\bar H_{k+1}(B')\bar\Psi'$,
\begin{equation}\label{eq:rg-step-induced-fields}
 H_k(B)\Psi^c=\psi'_L,\qquad
 \bar H_k(B)\bar\Psi^c=\bar\psi'_L.
\end{equation}
The first relation is \eqref{eq:rg-step-coarse-scaling}.
For the second use \eqref{eq:rg-step-transpose}
and apply the first relation at $-B$ to
$\mathsf C\bar\Psi'$; the constant matrix
$\mathsf C$ commutes with scaling and the averages.
The covariant H\"older quotients of these induced
fields are also the scaled outgoing quotients:
the transports agree and $d_X$ is divided by $L$,
so the quotient carries the factor $L^{-1-\alpha}$.
This holds for every pair, including a pair inside
one cluster.
Applying $T_kH_k=1$ and its barred version shows
that the means in \eqref{eq:rg-step-mean-fields}
are also $T_k(B)\psi'_L$ and $\bar T_k(B)\bar\psi'_L$.
At level zero these maps are identities on
the chamber field.

The quadratic term in the Gaussian block formula is
therefore exactly
\[
 \langle\bar\Psi'_L,D_{k+1}^0(B)\Psi'_L\rangle_L
                         =\langle\bar\Psi',D_{k+1}(B')\Psi'\rangle_1,
\]
since the pairings of scaled fields multiply by $L$.
Together with \exteq{73}{FermionBlockManuscript},
this identifies the entire outgoing quadratic density.

\begin{proposition}\label{prop:rg-step-determinant-normalization}
On the invertibility domain just specified,
\begin{equation}\label{eq:rg-step-determinant-normalization}
 \det D_{k+1}(B')=
           b^{4n_{k+1}}\frac{\det\mathcal D_k(B)}{\delta Z_k(B)},
 \qquad n_k=\ell^3,\quad n_{k+1}=(\ell/L)^3.
\end{equation}
In particular, with $c_k=b^{-4n_{k+1}}$,
\begin{equation}\label{eq:rg-step-zero-normalization}
 \frac{c_k}{\det D_k(0)}
              =\frac{1}{\det D_{k+1}(0)\delta Z_k(0)}.
\end{equation}
\end{proposition}

\begin{proof}
Apply Sylvester's determinant identity to
$b_L\mathcal A_k\Gamma_k$ and $\mathcal A_k^\sharp$
on the two different spaces. It gives
\[
 b_L^{-4n_{k+1}}\det D_{k+1}^0(B)
       =\det(1-b_L\mathcal A_k^\sharp\mathcal A_k\Gamma_k)
       =\det(\mathcal D_k\Gamma_k)
       =\frac{\det\mathcal D_k(B)}{\delta Z_k(B)}.
\]
The incoming determinant at $k=0$ has size $192n_0$;
the exponent of $b_L$ is the outgoing dimension
$4n_1$, as in every other step.
Equation \eqref{eq:rg-step-coarse-scaling} gives
$\det D_{k+1}^0(B)=L^{-4n_{k+1}}\det D_{k+1}(B')$.
Using $Lb_L=b$ proves
\eqref{eq:rg-step-determinant-normalization}.
At the zero connection all denominators are nonzero
by \cref{prop:rg-step-real-inverses}, so
\eqref{eq:rg-step-zero-normalization} follows.
\end{proof}

\subsection{The density with an external connection}
\label{subsec:rg-step-density}

For this subsection assume $0<\mu\le1/L$ and
$e_{k+1}=L^{1/2}e_k\le e_{\mathrm{prim}}$, with the hypotheses
of Propositions~\extref{3.3}{FermionBlockManuscript},
\extref{4.2}{FermionBlockManuscript} and Theorem
\extref{B}{MaxwellBlockManuscript}. We also impose
\exteq{43}{MaxwellBlockManuscript} and
\eqref{eq:rg-domain-charge-conditions}.
The inverse domain $\widetilde G_k$ is the domain
$\widetilde{\mathcal G}_k$ of the propagator estimates.
Let $A'\in\mathrm{Dom}_{k+1}$ be a real block field,
with arbitrary harmonic part, and let
$\mathscr C'\in\mathfrak X_{k+1}$.
Put
\[
 \mathscr A'=H_{k+1}^{\mathrm{ph}}A',\qquad
 \mathscr A=\mathscr A'_L,\qquad
 \mathscr C=\mathscr C'_L,\qquad
 \mathscr T'=\mathscr A'+\mathscr C',\qquad
 \mathscr T=\mathscr T'_L.
\]
The photon block objects are
\begin{equation}\label{eq:rg-step-photon-fields}
 \begin{split}
 A_k(\zeta)&=A_k^{\min}(A')+\mathrm{Sq}_k\zeta,\\
 \mathscr Z(\zeta)&=H_k^{\mathrm{ph}}\mathrm{Sq}_k\zeta,\qquad
 \mathscr B_k=H_k^{\mathrm{ph}}A_k^{\min}(A')
                                      =\mathscr A-d\omega_k.
 \end{split}
\end{equation}
Here $\omega_k$ is real and periodic, as in
\exteq{30}{MaxwellBlockManuscript}.
The truncated photon measure has the normalization
\begin{equation}\label{eq:rg-step-box-normalization}
 \mu_I(\mathrm{Box}_k)=\exp(-n_k\delta\varepsilon_k^0),\qquad
 \widehat\mu_k=
              \frac{\mathbf1_{\mathrm{Box}_k}}{\mu_I(\mathrm{Box}_k)}\mu_I,
\end{equation}
proved in \exteq{24}{MaxwellBlockManuscript}.
Its denominator is positive: $e_k<1$ makes
$p_{0,k}=(-\log e_k)^{p_0}>0$, and the box has
nonempty interior in the photon coordinate space.

All connections read by the step have their required
domains. Indeed, scaling gives
$\mathscr A\in\frac14L^{-3/4}\widetilde R_k
\subset\widetilde R_k^-$ and
$\mathscr C\in\mathfrak X_k$.
Theorem~\extref{B}{MaxwellBlockManuscript}
gives $A_k(\zeta)\in\mathrm{Dom}_k$ for every
$\zeta\in\mathrm{Box}_k$.
The real-gauge identity in
\eqref{eq:rg-step-photon-fields} yields
\[
 H_k^{\mathrm{ph}}A_k(\zeta)
                 =\mathscr A+\mathscr Z(\zeta)-d\omega_k.
\]
Since $\frac14\widetilde R_k$ is invariant under
real gauges, $\mathscr A+\mathscr Z(\zeta)$ also
lies in this domain. Thus both connection pairs
$(\mathscr A,\mathscr C)$ and
$(\mathscr A+\mathscr Z,\mathscr C)$ lie in
$\widetilde R_k\times\mathfrak X_k$.
By \eqref{eq:rg-total-domain-inclusion},
$\mathscr T$, $\mathscr T+\mathscr Z$ and
$\mathscr T+\mathscr Z-d\omega_k$, and their negatives,
belong to $\widetilde G_k$.
The fine and coarse inverse estimates consequently
apply to them. The outgoing operators at
$\mathscr T'$ then exist by
\eqref{eq:rg-step-fine-scaling}.

For a total connection $B$ let $\operatorname{ev}_B$
insert
\[
 \psi=H_k(B)\Psi,\qquad
 \bar\psi=\bar H_k(B)\bar\Psi,\qquad
 \chi=\delta_{\alpha,B}H_k(B)\Psi,\qquad
 \bar\chi=\delta_{\alpha,-B}\bar H_k(B)\bar\Psi.
\]
At $k=0$ the two field maps are identities.
For complex $m,\varepsilon^{\mathrm v}$ and
$\mathscr E\in K_k^{x,\C}$, with $\mathscr E=0$
at $k=0$, define
\begin{equation}\label{eq:rg-step-interaction-action}
 \begin{split}
 I_k^x[m,\mathscr E](A,C;\Psi)
     &=\operatorname{ev}_{A+C}
         \left(-m\langle\bar\psi,\psi\rangle
                               +\sum_X\mathscr E(X;A,C)\right),\\
 \mathrm{Act}_k^x[m,\mathscr E](A,C;\Psi)
     &=-\langle\bar\Psi,D_k(A+C)\Psi\rangle
                                  +I_k^x[m,\mathscr E](A,C;\Psi).
 \end{split}
\end{equation}
The mass pairing is always the chamber pairing
$\eta_k^3\,48^{-1}\sum_{c,g,\beta}\bar\psi_\beta(c,g)\psi_\beta(c,g)$.
The kinetic pairing uses the measure of $\mathscr V_k$.
For an arbitrary finite family of scalar functions
$\mathfrak P_k(X,C)$ on $\mathfrak X_k$, set
\begin{equation}\label{eq:rg-step-input-density}
 \begin{split}
 \mathscr F_k^x[\varepsilon^{\mathrm v},m,\mathscr E](A,C;\Psi)
   =\det D_k(0)^{-1}\exp\left\{
         \mathrm{Act}_k^x[m,\mathscr E](H_k^{\mathrm{ph}}A,C;\Psi)
                   -\varepsilon^{\mathrm v}n_k
                              +\sum_X\mathfrak P_k(X,C)\right\}.
 \end{split}
\end{equation}
The arguments here are $A\in\mathrm{Dom}_k$ and
$C\in\mathfrak X_k$. The fields $\Psi$ and $\bar\Psi$
are independent Grassmann variables.

Define the block transformation on an even density
$F$ whose coefficients are integrable along
$A_k(\zeta)$ by
\begin{equation}\label{eq:rg-step-block-transformation}
 \begin{split}
 (\mathbf T_k^xF)(A',\mathscr C';\Psi')
  &=c_k\int_{\mathrm{Box}_k}\dd\mu_I(\zeta)
                       \int\mathrm D\Psi\,e^{\mathcal Q_{\mathrm{bl}}}
                                         F(A_k(\zeta),\mathscr C;\Psi),\\
 \mathcal Q_{\mathrm{bl}}
  &=-b_L\big\langle
       e^{-ie_k\omega_k}\bar\Psi'_L
                    -\mathcal A_k(-\mathscr B_k-\mathscr C)\bar\Psi,\,
       e^{ie_k\omega_k}\Psi'_L
                    -\mathcal A_k(\mathscr B_k+\mathscr C)\Psi
                                                \big\rangle_L.
 \end{split}
\end{equation}
The gauge function in the two outgoing phases is
evaluated at the $L$-spaced sites.
The Berezin functional is the weighted functional
of \exteq{65}{FermionBlockManuscript}; its incoming mode
measure is $1$ for $k\ge1$ and $1/48$ for $k=0$.

Its normalization is
\begin{equation}\label{eq:rg-step-block-marginal}
 \int\mathrm D\Psi'\,(\mathbf T_k^xF)(A',\mathscr C';\Psi')
       =\int_{\mathrm{Box}_k}\dd\mu_I(\zeta)
                       \int\mathrm D\Psi\,F(A_k(\zeta),\mathscr C;\Psi).
\end{equation}
To prove this, integrate the outgoing family first.
Reciprocal multiplication of its barred and unbarred
generators removes the two phases and leaves the top
coefficient unchanged. Translation by the two averaged
incoming fields likewise leaves it unchanged: every
translated term that contains an external odd field
has lower degree in the outgoing family.
The remaining outgoing Gaussian exponent is
$-b_L\langle\bar\Psi'_L,\Psi'_L\rangle_L
=-b\langle\bar\Psi',\Psi'\rangle_1$.
Its integral is $b^{4n_{k+1}}=c_k^{-1}$ by
\exteq{68}{FermionBlockManuscript}.
The two Berezin integrations commute, since both
extract an even top-degree coefficient; integration
of the finite-dimensional coefficient functions in
$\zeta$ also commutes with them. This proves
\eqref{eq:rg-step-block-marginal}.

For the density \eqref{eq:rg-step-input-density}
the coefficient integrability follows from the domain
statements above. The bond entries of all inverse
operators are analytic on their invertibility domains,
and the activity kernels are holomorphic. Composing
them with the affine map $\zeta\mapsto
H_k^{\mathrm{ph}}A_k(\zeta)$ gives continuous
coefficients on the compact photon box. Their even
Grassmann exponentials are also continuous there.

\begin{proposition}\label{prop:rg-step-density-identity}
Let the data and domains be as above. Put
\[
 m'=m+m(\mathscr E),\qquad \mathscr E^r=R\mathscr E,\qquad
 m^+=Lm',\qquad
 \varepsilon^{\mathrm v,+}
       =L^3\big(\varepsilon^{\mathrm v}
              +\varepsilon_{\mathrm v}(\mathscr E)+\delta\varepsilon_k^0\big).
\]
At $k=0$ the activity and its extracted terms vanish.
Define the outgoing inserted fields
$\psi'=H_{k+1}(\mathscr T')\Psi'$,
$\bar\psi'=\bar H_{k+1}(\mathscr T')\bar\Psi'$ and
their covariant H\"older quotients, and let
$\psi^{0,\#}$ denote their scaled fine fields,
including the barred ones.
Let the means $\Psi^c,\bar\Psi^c$ be those of
\eqref{eq:rg-step-mean-fields} at $B=\mathscr T$,
and put
\[
 \Psi^f=\Psi^c+\Gamma_k(\mathscr T)W,\qquad
 \bar\Psi^f=\bar\Psi^c+\bar W.
\]
With $\mathscr Z=H_k^{\mathrm{ph}}\mathrm{Sq}_k\zeta$,
set
\begin{equation}\label{eq:rg-step-fluctuation-exponent}
 \begin{split}
 V_k^x={}&
       \mathrm{Act}_k^x[m',\mathscr E^r]
                    (\mathscr A+\mathscr Z,\mathscr C;\Psi^f)
       -\mathrm{Act}_k^x[m',\mathscr E^r]
                    (\mathscr A,\mathscr C;\Psi^f)\\
       &+I_k^x[m',\mathscr E^r](\mathscr A,\mathscr C;\Psi^f)
       -I_k^x[m',\mathscr E^r](\mathscr A,\mathscr C;\Psi^c),\\
 \Xi_k^x={}&\int\dd\mu_I(W)\,\dd\widehat\mu_k(\zeta)\,e^{V_k^x}.
 \end{split}
\end{equation}
Then the exact identity is
\begin{equation}\label{eq:rg-step-density-identity}
 \begin{split}
 &(\mathbf T_k^x\mathscr F_k^x
                [\varepsilon^{\mathrm v},m,\mathscr E])
                                      (A',\mathscr C';\Psi')\\
 &\quad=\det D_{k+1}(0)^{-1}
       \exp\Big\{-\langle\bar\Psi',D_{k+1}(\mathscr T')\Psi'\rangle_1
                  -m^+\langle\bar\psi',\psi'\rangle
                  -\varepsilon^{\mathrm v,+}n_{k+1}\\
 &\hspace{48mm}
          +\sum_X\mathscr E^r(X;\mathscr A,\mathscr C;\psi^{0,\#})
          +\sum_X\mathfrak P_k(X,\mathscr C)\Big\}
                     \frac{\delta Z_k(\mathscr T)}{\delta Z_k(0)}\Xi_k^x.
 \end{split}
\end{equation}
All integrals in this formula are finite coefficient
integrals. Its right side depends on $A'$ through
$H_{k+1}^{\mathrm{ph}}A'$.
\end{proposition}

\begin{proof}
Extraction gives
\[
 I_k^x[m,\mathscr E](A,C;\Psi)
      =-\varepsilon_{\mathrm v}(\mathscr E)n_k
                                      +I_k^x[m',\mathscr E^r](A,C;\Psi).
\]
This is the image of
\exteq{18}{NormalizedContraction} under the
field evaluation homomorphism. At $k=0$ it holds
with zero activity and $m'=m$.

For a fixed $\zeta$, change the integrated variables
by $\Psi\mapsto e^{ie_k\omega_k}\Psi$ and
$\bar\Psi\mapsto e^{-ie_k\omega_k}\bar\Psi$.
This leaves the weighted Berezin functional unchanged,
since each mode's factors cancel in its top monomial,
also for the $192n_0$ modes of weight $1/48$.
The transport identity
\[
 U_{B-d\omega_k}(c,c')e^{ie_k\omega_k(c')}
                       =e^{ie_k\omega_k(c)}U_B(c,c')
\]
intertwines the Dirac operator and every average.
Products and inverses preserve this intertwining,
so it also holds for $H_k$ and $D_k$.
Its covariant difference identity uses the phase
at the first photon site, including for a pair
inside one cluster. The barred identity has the
inverse phase. Gauge covariance of the activity
and invariance of the mass pairing consequently
turn its first connection argument
$\mathscr A+\mathscr Z-d\omega_k$ into
$\mathscr A+\mathscr Z$.
The same change removes the gauge from the block
Gaussian factor. The integral in
\eqref{eq:rg-step-block-transformation} is now
\[
 \int\mathrm D\Psi\,
       \exp\left\{-b_L
          \langle\bar\Psi'_L-\mathcal A_k(-\mathscr T)\bar\Psi,\,
                  \Psi'_L-\mathcal A_k(\mathscr T)\Psi\rangle_L
            +\mathrm{Act}_k^x[m',\mathscr E^r]
                    (\mathscr A+\mathscr Z,\mathscr C;\Psi)\right\},
\]
multiplied by
$\det D_k(0)^{-1}
\exp\{-(\varepsilon^{\mathrm v}+
\varepsilon_{\mathrm v}(\mathscr E))n_k+
\sum_X\mathfrak P_k(X,\mathscr C)\}$.

Separate the background quadratic term by writing
\[
 \begin{split}
 \mathcal R(\Psi)={}&I_k^x[m',\mathscr E^r](\mathscr A,\mathscr C;\Psi)\\
  &+\mathrm{Act}_k^x[m',\mathscr E^r]
                    (\mathscr A+\mathscr Z,\mathscr C;\Psi)
   -\mathrm{Act}_k^x[m',\mathscr E^r]
                    (\mathscr A,\mathscr C;\Psi).
 \end{split}
\]
The last exponent above is then the background
Gaussian exponent plus $\mathcal R(\Psi)$.
Translation by the means \eqref{eq:rg-step-mean-fields}
and the block formula \exteq{71}{FermionBlockManuscript}
give
\[
 \delta Z_k(\mathscr T)
   e^{-\langle\bar\Psi',D_{k+1}(\mathscr T')\Psi'\rangle_1}
          \int\dd\mu_I(W)\,e^{\mathcal R(\Psi^f)}.
\]
The change of variables in this formula uses the
same mode measure for $\Psi$ and $W$.
At $k=0$ it is the identity
\[
 \int\mathrm D\Psi\,e^{-\langle\bar\Psi,F_0(\mathscr T)\Psi\rangle^{\mathrm{ch}}}
       G(\bar\Psi,\Psi)
   =\det F_0(\mathscr T)
          \int\dd\mu_I(W)\,
                    G(\bar W,\Gamma_0(\mathscr T)W),
\]
on the entire $192n_0$-mode space.
The Gaussian measure is
$e^{-\langle\bar W,W\rangle^{\mathrm{ch}}}\mathrm DW$,
with the weight $\sqrt{48}$ in its coefficient bound.

By definition of $V_k^x$,
$\mathcal R(\Psi^f)=
I_k^x[m',\mathscr E^r](\mathscr A,\mathscr C;\Psi^c)+V_k^x$.
Its first summand is independent of $W$ and $\zeta$,
so its even exponential factors out.
The induced-field identity
\eqref{eq:rg-step-induced-fields} gives
\[
 I_k^x[m',\mathscr E^r](\mathscr A,\mathscr C;\Psi^c)
      =-m'\langle\bar\psi'_L,\psi'_L\rangle
                 +\sum_X\mathscr E^r(X;\mathscr A,\mathscr C;\psi^{0,\#}).
\]
The chamber measure changes by $L^3$ and each field
by $L^{-1}$. Hence
$\langle\bar\psi'_L,\psi'_L\rangle
=L\langle\bar\psi',\psi'\rangle$, giving $m^+=Lm'$.
Normalizing the photon-box restriction in
\eqref{eq:rg-step-box-normalization} adds
$-\delta\varepsilon_k^0 n_k$ to the vacuum term.
Since $n_k=L^3n_{k+1}$, this term has exactly
the stated outgoing vacuum coefficient.
Finally \eqref{eq:rg-step-zero-normalization}
turns the remaining constant factor into
$\det D_{k+1}(0)^{-1}\delta Z_k(0)^{-1}$,
proving \eqref{eq:rg-step-density-identity}.

The fluctuation exponent is even and continuous
on the photon box, by the same coefficient argument
as for the input density. Thus its integral is defined,
and its two finite integrations commute.
Every background in the final expression is formed
from $\mathscr A'$ and $\mathscr C'$; the remaining
photon maps and measure are independent of $A'$.
The minimum and the real gauge have disappeared
from that expression, proving the last assertion.
\end{proof}

\subsection{Local substitution of the mean field}
\label{subsec:rg-step-local-substitution}

We recall the result of \cite[Theorem~B(iii)--(iv)]{NormalizedContraction}
that is used between the two localizations. For $k\ge1$ and unit
generators over $Y$, the substitution fixes the old ordinary and
pair generators and sets
\begin{equation}\label{eq:rg-step-local-substitution}
 \operatorname{sub}_{Y,B}\Psi=T_Y^{(1)}(B)\psi+T_Y^{(2)}(B)\chi,
 \qquad
 \operatorname{sub}_{Y,B}\bar\Psi=
              \bar T_Y^{(1)}(B)\bar\psi+\bar T_Y^{(2)}(B)\bar\chi.
\end{equation}
The row measures are $\eta_k^3/48$ for ordinary columns and
$(\eta_k^3/48)^2$ for pair columns. The rows add only ordinary
legs over $Y$ and boundary pairs $(x,x^{(3)})$ admissible for $Y$;
they read only bonds meeting $Y$. The substitution preserves
balance, parity and the lattice, local gauge, charge and reality
identities of the old generators. In particular,
\begin{equation}\label{eq:rg-step-substitution-covariances}
 \begin{split}
 a_g\operatorname{sub}_{gY,gB}
                        &=\operatorname{sub}_{Y,B}a_g,\\
 G_\lambda\operatorname{sub}_{Y,B^1}
                        &=\operatorname{sub}_{Y,B}G_\lambda,
                 \qquad B^1=B-d\lambda\text{ on bonds meeting }Y,\\
 \mathrm{Cg}\operatorname{sub}_{Y,-B}
                        &=\operatorname{sub}_{Y,B}\mathrm{Cg},\\
 \mathbf K\operatorname{sub}_{Y,-\overline B}
                        &=\operatorname{sub}_{Y,B}\mathbf K.
 \end{split}
\end{equation}
The local gauge equality need hold only on the bonds meeting $Y$.
For $B\in\widetilde R_k+\mathfrak X_k$, the row bound gives
\begin{equation}\label{eq:rg-step-substitution-bound}
 \|\operatorname{sub}_{Y,B}F\|_{w_1,w_2}
       \le\|F\|_{w_1,w_2;w_3}
       \quad\text{if }C_T(w_1+w_2)\le w_3.
\end{equation}
The bounds on imaginary part and curvature in
\eqref{eq:rg-total-imaginary-curvature-bounds} verify the hypotheses
of that result. The primitive inverse domain at $B$ and $-B$
also verifies its additional recovery hypothesis. Therefore, on
evaluation at $\psi=H_k(B)\Phi$, $\chi=\delta_{\alpha,B}\psi$ and
the barred fields, the substitution sends $\Psi,\bar\Psi$ to
$\Phi,\bar\Phi$. At zero no such substitution is needed.

